% Generated by revision/make_paper_tables.py -- do not edit by hand.
% Regenerate after re-running the Phase 1 task scripts.
\begin{table}[H]
\centering
\begin{threeparttable}
\caption{The comparison at three levels of stratification. The absolute level falls as
the locus effect is removed; the margin grows.}
\label{tab:stratified}
\footnotesize
\begin{tabular}{l S[table-format=3.0] S[table-format=1.4] S[table-format=1.4] S[table-format=+1.4] c}
\toprule
Evaluation & {Strata} & {OptiPrime} & {PE-RankFormer} & {$\Delta\rhoS$} & {95\% CI} \\
\midrule
Pooled (Table~\ref{tab:main}) & 1 & 0.8690 & 0.9079 & +0.0389 & $[+0.0288,\,+0.0498]$ \\
Within source study, $n$-weighted & 2 & 0.7788 & 0.8331 & +0.0543 & {---} \\
Within condition, $n$-weighted & 14 & 0.7605 & 0.8111 & +0.0506 & $[+0.0345,\,+0.0670]$ \\
\textbf{Within target}, mean over targets & 670 & \bfseries 0.5472 & \bfseries 0.6356 & \bfseries +0.0884 & $[+0.0761,\,+0.1007]$ \\
\bottomrule
\end{tabular}
\begin{tablenotes}[flushleft]\footnotesize
\item A condition is a (source study, cell line, PE system) cell; a target is a
protospacer. A target is scored if it carries at least 5 designs and a non-constant
measured outcome: 670 of 750 protospacers qualify, covering 19,334 of the
20{,}509 held-out rows, with a median of 18 designs per target (range 6--200).
\item The within-target row is an unweighted mean over targets, so each design problem
counts once, which is the deployment-relevant weighting; weighting by target size
instead gives +0.0758 (95\% CI $[+0.0649,\,+0.0874]$).
\item Intervals are paired bootstraps over 5{,}000 resamples of the 750 protospacer
clusters. PE-RankFormer leads on 78.4\% of individual targets, in all 14
conditions, and in 100\% of resamples at both the condition and the target level.
\item Absolute levels fall for \emph{both} models as stratification tightens, because
between-stratum differences in efficiency level stop contributing to the correlation.
The within-target row is the quantity relevant to prospective design.
\item Reproduce with \path{revision/task_1_1_per_target.py}.
\end{tablenotes}
\end{threeparttable}
\end{table}
